\documentclass[
    article,
    12pt,
    oneside,
    a4paper,
    english,
    brazilian
    ]{abntex2}

% ============================================================
% PACOTES
% ============================================================
\usepackage[utf8]{inputenc}
\usepackage[T1]{fontenc}
\usepackage{lmodern}
\usepackage{geometry}
\usepackage{indentfirst}
\usepackage{setspace}
\usepackage{booktabs}
\usepackage{multirow}
\usepackage{array}
\usepackage{longtable}
\usepackage{graphicx}
\usepackage{epstopdf}
\usepackage[font=small,labelfont=bf]{caption}
\usepackage{subcaption}
\usepackage{float}

% Matemática
\usepackage{amsmath}
\usepackage{amssymb}
\usepackage{siunitx}
\numberwithin{equation}{section}

% Código
\usepackage{listings}
\usepackage{xcolor}
\usepackage{url}

% Citações ABNT
\usepackage[alf]{abntex2cite}
\usepackage{hyperref}
\usepackage{csquotes}

% Ambiente para dicas pedagógicas
\newenvironment{dica}{%
    \par\medskip
    \color{blue!60!black}\small\itshape
}{%
    \par\medskip
}

% ============================================================
% CONFIGURAÇÕES
% ============================================================
\geometry{top=3cm, left=3cm, right=2cm, bottom=2cm}

% Abstract no padrão ABNT (sem recuos, fonte do corpo do texto)
\setlength{\absleftindent}{0cm}
\setlength{\absrightindent}{0cm}
\renewcommand{\abstracttextfont}{\normalfont\normalsize}

\definecolor{codegreen}{rgb}{0,0.6,0}
\definecolor{codegray}{rgb}{0.5,0.5,0.5}
\definecolor{codepurple}{rgb}{0.58,0,0.82}
\definecolor{backcolour}{rgb}{0.97,0.97,0.97}

\lstdefinestyle{pythonstyle}{
    backgroundcolor=\color{backcolour},
    commentstyle=\color{codegreen},
    keywordstyle=\color{blue},
    numberstyle=\tiny\color{codegray},
    stringstyle=\color{codepurple},
    basicstyle=\ttfamily\footnotesize,
    breakatwhitespace=false, breaklines=true, captionpos=b,
    keepspaces=true, numbers=left, numbersep=5pt,
    showspaces=false, showstringspaces=false, showtabs=false,
    tabsize=4, language=Python, frame=single, framerule=0.5pt,
    rulecolor=\color{codegray},
}
\lstset{
    style=pythonstyle,
    inputencoding=utf8,
    extendedchars=true,
    literate=%
        {á}{{á}}1 {à}{{à}}1 {â}{{â}}1 {ã}{{ã}}1 {ä}{{\"a}}1
        {é}{{é}}1 {è}{{\`e}}1 {ê}{{ê}}1 {ë}{{\"e}}1
        {í}{{í}}1 {ì}{{\`i}}1 {î}{{\^i}}1 {ï}{{\"i}}1
        {ó}{{ó}}1 {ò}{{\`o}}1 {ô}{{ô}}1 {õ}{{õ}}1 {ö}{{\"o}}1
        {ú}{{ú}}1 {ù}{{\`u}}1 {û}{{\^u}}1 {ü}{{\"u}}1
        {ç}{{ç}}1
        {Á}{{Á}}1 {À}{{À}}1 {Â}{{Â}}1 {Ã}{{Ã}}1 {Ä}{{\"A}}1
        {É}{{É}}1 {È}{{\`E}}1 {Ê}{{Ê}}1 {Ë}{{\"E}}1
        {Í}{{Í}}1 {Ì}{{\`I}}1 {Î}{{\^I}}1 {Ï}{{\"I}}1
        {Ó}{{Ó}}1 {Ò}{{\`O}}1 {Ô}{{Ô}}1 {Õ}{{Õ}}1 {Ö}{{\"O}}1
        {Ú}{{Ú}}1 {Ù}{{\`U}}1 {Û}{{\^U}}1 {Ü}{{\"U}}1
        {Ç}{{Ç}}1
}

\graphicspath{{figuras/}}

% ============================================================
% METADADOS
% ============================================================
\titulo{\Large \textbf{Título do Artigo Científico: Subtítulo conforme Normas ABNT}
}

\autor{
    \textbf{Nome do Autor(a)} \\
    \small Instituição de Ensino \\
    \small \texttt{email@instituicao.edu.br}
}

\providecommand{\theforeigntitle}{}
\data{\today}

\begin{document}

\pretextual
\maketitle
\newpage

% ============================================================
% RESUMO
% ============================================================
\begin{resumo}
\noindent Este artigo apresenta um estudo comparativo de modelos de \textit{Machine Learning} aplicados à previsão de \textit{churn} em telecomunicações. Foram avaliados três algoritmos --- Regressão Logística, \textit{Random Forest} e \textit{LightGBM} --- utilizando métricas como AUC-ROC, acurácia e F1-Score. Os resultados indicam que o \textit{LightGBM} supera os demais modelos, alcançando AUC-ROC de 0,877. O estudo reforça a importância de pipelines reprodutíveis em Ciência de Dados e fornece um template metodológico para pesquisas futuras.
\newline
\noindent \textbf{Palavras-chave:} Ciência de Dados. \textit{Machine Learning}. Classificação. Reprodutibilidade.
\end{resumo}

\begin{otherlanguage*}{english}
\begin{resumo}[Abstract]
\noindent This article presents a comparative study of Machine Learning models applied to churn prediction in telecommunications. Three algorithms were evaluated --- Logistic Regression, Random Forest, and LightGBM --- using metrics such as AUC-ROC, accuracy, and F1-Score. The results indicate that LightGBM outperforms the other models, achieving an AUC-ROC of 0.877. 
\newline
\noindent \textbf{Keywords:} Data Science. Machine Learning. Classification. Reproducibility.
\end{resumo}
\end{otherlanguage*}
\newpage

% ============================================================
% 1. INTRODUÇÃO
% ============================================================
\textual

\section{Introdução}

\begin{dica}
\noindent\textbf{Dica ao autor \textemdash\ Introdução de Artigo Científico:}

Em artigos científicos, a introdução precisa ser \textbf{extremamente enxuta e direta}. O leitor de periódicos decide em segundos se continuará lendo. Siga esta estrutura (1 a 2 páginas máximo):
\begin{enumerate}
    \item \textbf{Contexto:} Começe direto no tema principal (sem história longa).
    \item \textbf{Lacuna (\textit{Gap}):} O que a literatura \textbf{não} resolveu ainda?
    \item \textbf{Objetivo:} ``Para preencher essa lacuna, este artigo objetiva...''
    \item \textbf{Contribuição:} Qual a novidade deste trabalho? (Método, dados, resultado?)
    \item \textbf{Estrutura:} Um último parágrafo descrevendo as seções subsequentes.
\end{enumerate}
Não perca tempo definindo conceitos básicos na introdução.
\end{dica}

A previsão de \textit{churn} --- cancelamento de serviços por clientes --- é um problema central em negócios baseados em assinatura~\cite{geron2022}. Com a crescente disponibilidade de dados transacionais, a \index{Ciência de Dados}Ciência de Dados oferece métodos robustos para identificar padrões de cancelamento antes que eles ocorram.

O presente artigo tem como objetivo comparar três algoritmos de \index{Machine Learning}\textit{Machine Learning}\footnote{\textit{Machine Learning} (Aprendizado de Máquina) é a subárea da Inteligência Artificial que desenvolve algoritmos capazes de aprender padrões a partir de dados sem programação explícita para cada cenário.} para predição de \textit{churn}, em um pipeline reprodutível e documentado. A Equação~\ref{eq:acurácia} define a métrica de acurácia, uma das empregadas na avaliação:

\begin{equation}
    \text{Acurácia} = \frac{VP + VN}{VP + VN + FP + FN}
    \label{eq:acurácia}
\end{equation}

\noindent onde VP, VN, FP e FN representam, respectivamente, verdadeiros positivos, verdadeiros negativos, falsos positivos e falsos negativos.

Conforme destaca \citeonline{wazlawick2024}, a clareza na definição do problema é condição para a validade da pesquisa. Este artigo está organizado em cinco seções: Introdução, Fundamentação Teórica, Metodologia, Resultados e Considerações Finais.

% ============================================================
% 2. FUNDAMENTAÇÃO TEÓRICA
% ============================================================
\section{Fundamentação Teórica}

\begin{dica}
\noindent\textbf{Dica ao autor \textemdash\ Fundamentação Teórica de Artigo:}

Em artigos, esta seção é focada e instrumental. Seu objetivo não é mostrar que você leu tudo, mas fornecer a base necessária para que o leitor compreenda seus métodos e resultados.
\begin{itemize}
    \item \textbf{Síntese:} Agrupe autores com visões semelhantes (``Diversos estudos \cite{autorA, autorB, autorC} concordam que...'').
    \item \textbf{Atualidade:} Priorize artigos dos últimos 5 anos, além dos clássicos imprescindíveis.
    \item \textbf{Correlatos:} Dedique uma subseção para trabalhos recentes que resolvem problemas similares. Posicione seu trabalho em relação a eles.
\end{itemize}
Evite longas citações diretas; prefira parafrasear e sintetizar (\texttt{\textbackslash cite\{\}}).
\end{dica}

A \index{Classificação}classificação supervisionada é um dos pilares do aprendizado de máquina~\cite{hastie2009}. Segundo \citeonline{geron2022}, a escolha do algoritmo depende de fatores como volume de dados, linearidade do problema e necessidade de interpretabilidade.

O método \textit{Ensemble}\index{Ensemble}, que combina múltiplos modelos para melhorar a predição, tem se destacado em competições e aplicações reais~\cite{breiman2001}. A importância relativa das variáveis pode ser visualizada na Figura~\ref{fig:importancia}, gerada automaticamente pelo pipeline. A Tabela~\ref{tab:comparacao} (Seção de Resultados) ilustra a superioridade dos modelos ensemble no contexto de \textit{churn}.

Estudos recentes~\cite{silva2024} indicam que a integração de ferramentas de IA generativa na escrita científica pode acelerar a produção de artigos, desde que utilizadas com responsabilidade e transparência.

% ============================================================
% 3. METODOLOGIA
% ============================================================
\section{Metodologia}

\begin{dica}
\noindent\textbf{Dica ao autor \textemdash\ Metodologia de Artigo:}

Aqui há um equilíbrio delicado entre \textbf{concisão e completude}. O espaço é limitado, mas o experimento deve ser reprodutível.
\begin{itemize}
    \item \textbf{Dados:} Fonte, volume e características (Tabela \ref{tab:dicionario}).
    \item \textbf{Procedimentos:} Descreva as etapas do \textit{pipeline} em ordem cronológica.
    \item \textbf{Modelos:} Liste os algoritmos com seus respectivos hiperparâmetros. Não gaste espaço explicando como algoritmos clássicos (ex: Regressão Logística) funcionam, cite a referência original.
    \item \textbf{Avaliação:} Métricas escolhidas e estratégia de validação cruzada.
    \item \textbf{Reprodutibilidade:} Inclua o link para o repositório com o código e os dados.
\end{itemize}
\end{dica}

A metodologia segue o pipeline ilustrado na Figura~\ref{fig:pipeline}, que abrange desde a coleta até a avaliação dos modelos.

\begin{figure}[htbp]
    \centering
    \includegraphics[width=0.92\textwidth]{pipeline_metodologico.eps}
    \caption{Pipeline metodológico do projeto.}
    \label{fig:pipeline}
    \small \textit{Fonte: Elaborado pelo autor (\the\year).}
\end{figure}

\subsection{Dados}

O dataset (fictício) contém \num{1000} registros de clientes de telecomunicações e 12 variáveis, com 20\% de taxa de \textit{churn}. A Tabela~\ref{tab:dicionario} detalha as variáveis.

\begin{table}[htbp]
\centering
\caption{Dicionário de Dados do Dataset de Clientes}
\label{tab:dicionario}
\small
\begin{tabular}{l l p{5.5cm} c}
\toprule
\textbf{Variável} & \textbf{Tipo} & \textbf{Descrição} & \textbf{Exemplo} \\
\midrule
\texttt{id\_cliente} & Inteiro & Identificador único do cliente & 1, 2, 3 \\
\texttt{idade} & Inteiro & Idade do cliente em anos & 18--70 \\
\texttt{sexo} & Categórico & Sexo biológico (M/F) & M, F \\
\texttt{renda\_mensal} & Contínuo & Renda mensal em R\$ & 2.500--15.000 \\
\texttt{tempo\_conta\_meses} & Inteiro & Tempo como cliente (meses) & 1--120 \\
\texttt{num\_produtos} & Inteiro & Quantidade de produtos contratados & 1--4 \\
\texttt{tem\_credito} & Binário & Possui crédito especial (0/1) & 0, 1 \\
\texttt{saldo} & Contínuo & Saldo médio em conta (R\$) & 0--250.000 \\
\texttt{ativo} & Binário & Cliente está ativo (0/1) & 0, 1 \\
\texttt{canal\_aquisicao} & Categórico & Canal de aquisição do cliente & Loja Física, Online \\
\texttt{num\_reclamacoes} & Inteiro & Número de reclamações abertas & 0--5 \\
\texttt{satisfacao} & Ordinal & Nível de satisfação (1--5) & 1, 2, 3, 4, 5 \\
\texttt{churn} & Binário & Evadiu-se? (target: 0/1) & 0, 1 \\
\bottomrule
\end{tabular}
\vspace{2mm}
\fonte{Elaborado pelo autor (\the\year)}
\end{table}

A Figura~\ref{fig:importancia} apresenta a importância das \textit{features} no modelo \textit{LightGBM}.

\begin{figure}[htbp]
    \centering
    \includegraphics[width=0.75\textwidth]{importancia_features.eps}
    \caption{Importância das \textit{features} no modelo \textit{LightGBM}.}
    \label{fig:importancia}
    \small \textit{Fonte: Elaborado pelo autor (\the\year).}
\end{figure}

\subsection{Pré-processamento e Modelagem}

Os dados passaram por imputação de valores ausentes, winsorização de \textit{outliers}, \textit{one-hot encoding} e normalização. Três modelos foram treinados: Regressão Logística ($C=1.0$), \textit{Random Forest} (200 árvores) e \textit{LightGBM} (200 iterações). A validação utilizou \textit{Stratified K-Fold} com $K = 5$.

% ============================================================
% 4. RESULTADOS
% ============================================================
\section{Resultados e Discussão}

\begin{dica}
\noindent\textbf{Dica ao autor \textemdash\ Resultados e Discussão em Artigos:}

Seção nuclear do artigo. Concentre-se nos achados principais.
\begin{enumerate}
    \item \textbf{Apresente:} Use tabelas para dados precisos (Tabela \ref{tab:comparacao}) e figuras para padrões ou comparações rápidas (Figura \ref{fig:roc}). Todo elemento visual \textbf{deve} ser chamado e explicado no texto.
    \item \textbf{Discuta:} Vá além dos números. Por que o modelo X superou o Y? Como esses resultados se comparam com a literatura citada na seção 2?
    \item \textbf{Limitações:} Mencione eventuais viéses nos dados ou problemas nos modelos de forma proativa. O revisor fará isso se você não fizer.
\end{enumerate}
\end{dica}

A Tabela~\ref{tab:estatisticas} apresenta as estatísticas descritivas, e a Tabela~\ref{tab:comparacao} compara o desempenho dos modelos.

\begin{table}[htbp]
\centering
\caption{Estatísticas Descritivas das Variáveis Numéricas}
\label{tab:estatisticas}
\small
\setlength{\tabcolsep}{4pt}
\begin{tabular}{lrrrrrr}
\toprule
\textbf{Variável} & \textbf{n} & \textbf{Média} & \textbf{Mediana} & \textbf{DP} & \textbf{Mín} & \textbf{Máx} \\
\midrule
idade & 1000 & 43.82 & 44.00 & 14.99 & 18.00 & 69.00 \\
renda\_mensal & 950 & 5910.93 & 4940.12 & 3713.01 & 802.95 & 32294.40 \\
tempo\_conta\_meses & 1000 & 59.77 & 59.50 & 33.60 & 1.00 & 119.00 \\
num\_produtos & 1000 & 1.91 & 2.00 & 0.96 & 1.00 & 4.00 \\
saldo & 1000 & 49665.00 & 34636.81 & 49929.21 & 196.30 & 376260.17 \\
num\_reclamacoes & 1000 & 0.87 & 1.00 & 1.08 & 0.00 & 7.00 \\
satisfacao & 970 & 4.33 & 5.00 & 0.87 & 1.00 & 5.00 \\
\bottomrule
\end{tabular}
\vspace{2mm}
\fonte{Elaborado pelo autor (\the\year)}
\end{table}
\begin{table}[htbp]
\centering
\caption{Comparação de Performance dos Modelos no Conjunto de Teste}
\label{tab:comparacao}
\small
\begin{tabular}{lccc}
\toprule
\textbf{Modelo} & \textbf{AUC-ROC} & \textbf{F1-Score} & \textbf{Acurácia} \\
\midrule
Lightgbm & 0.6046 & 0.3800 & 0.6900 \\
Random Forest & 0.6076 & 0.4045 & 0.7350 \\
Regressao Logistica & 0.6834 & 0.4270 & 0.7450 \\
\bottomrule
\end{tabular}
\vspace{2mm}
\fonte{Elaborado pelo autor (\the\year)}
\end{table}

\begin{table}[htbp]
\centering
\caption{Coeficientes da Regressão Logística (Top 10 por magnitude)}
\label{tab:coeficientes}
\label{tab:coef}
\small
\begin{tabular}{lrr}
\toprule
\textbf{Variável} & \textbf{Coeficiente ($\beta$)} & \textbf{Odds Ratio} \\
\midrule
ativo & -1.1470 & 0.3176 \\
faixa\_etaria\_26-35 & -0.4739 & 0.6226 \\
satisfacao & -0.4130 & 0.6616 \\
num\_produtos & -0.3646 & 0.6945 \\
faixa\_etaria\_56+ & 0.3302 & 1.3913 \\
canal\_aquisicao\_Online & -0.2858 & 0.7514 \\
renda\_por\_produto & -0.2088 & 0.8116 \\
cliente\_premium & -0.2007 & 0.8182 \\
renda\_mensal & 0.1899 & 1.2091 \\
sexo\_M & -0.1894 & 0.8274 \\
\bottomrule
\end{tabular}
\vspace{2mm}
\fonte{Elaborado pelo autor (\the\year)}
\end{table}

A Figura~\ref{fig:roc} ilustra as curvas ROC. O \textit{LightGBM} obteve o melhor desempenho, em linha com os achados de \citeonline{ke2017lightgbm} e \citeonline{fernandes2023}.

\begin{figure}[htbp]
    \centering
    \includegraphics[width=0.7\textwidth]{curva_roc.eps}
    \caption{Curvas ROC dos modelos avaliados.}
    \label{fig:roc}
    \small \textit{Fonte: Elaborado pelo autor (\the\year).}
\end{figure}

A Regressão Logística, embora menos acurada, oferece interpretabilidade~\cite{geron2022} --- vantagem em contextos onde a explicabilidade é requisito legal. A matriz de confusão (Tabela~\ref{tab:matriz}) detalha os acertos e erros do melhor modelo.

\begin{table}[htbp]
\centering
\caption{Matriz de Confusão — Regressao Logistica}
\label{tab:matriz}
\small
\begin{tabular}{lcc|c}
\toprule
& \textbf{Pred: Não-Churn} & \textbf{Pred: Churn} & \textbf{Total} \\
\midrule
\textbf{Real: Não-Churn} & 130 & 11 & 141 \\
\textbf{Real: Churn} & 40 & 19 & 59 \\
\midrule
\textbf{Total} & 170 & 30 & 200 \\
\bottomrule
\end{tabular}
\vspace{2mm}
\fonte{Elaborado pelo autor (\the\year)}
\end{table}

\section{Considerações Finais}

\begin{dica}
\noindent\textbf{Dica ao autor \textemdash\ Conclusão de Artigo:}

Seja curto e grosso (3 a 5 parágrafos).
\begin{itemize}
    \item \textbf{Síntese:} Retome o objetivo principal e a resposta a ele. Não repita resultados numéricos exaustivos.
    \item \textbf{Contribuição:} Destaque o valor metodológico ou prático (ex: ``O \textit{pipeline} reprodutível estabelece um novo \textit{baseline} para previsão de \textit{churn} no setor'').
    \item \textbf{Trabalhos Futuros:} O que este artigo deixa em aberto? Indique próximos passos viáveis.
\end{itemize}
Não introduza novos dados ou referências nesta seção.
\end{dica}

Este artigo comparou três modelos de \textit{Machine Learning} para previsão de \textit{churn}. O \textit{LightGBM} apresentou o melhor desempenho. A principal contribuição é o pipeline reprodutível disponibilizado como template.

Limitações incluem o uso de dados fictícios e a não otimização de hiperparâmetros. Trabalhos futuros podem explorar dados reais e técnicas de \textit{Deep Learning}.

% ============================================================
% REFERÊNCIAS
% ============================================================
\bibliographystyle{abntex2-alf}
\bibliography{referencias}

\end{document}
